\clearpage
\section*{September 9, 2026: the final fifteen construction reviews}
\addcontentsline{toc}{section}{September 9, 2026: the final fifteen construction reviews}
\textit{Agent-drafted entry.}

The fixed review cohort began this review with 253 closed cases and fifteen open cases. Research now covers all fifteen. The property-by-property evidence and recommendations are recorded in \texttt{tasks/working\_paper\_release\_audit/last\_fifteen\_review.md}. Research completion does not mean every recommendation has been applied or approved.

Five case resolutions have been entered. The ten-home building at 501--505 West 65th Place predates the study: City records show a 2000 new-building permit, passed inspections in 2002, and fire repairs in 2004. It is excluded despite a later Assessor year of 2006. Separate original permits and complete individual assessments support two three-apartment buildings at 6612 and 6616 South Michigan, retaining the Assessor's 2008 completion proxy. Their reported floor areas are 4,326 square feet each; their reported lots are 4,087 and 3,056 square feet respectively.

At 1002 and 1016 West Lake, completed residential assessments distinguish two apartments from one commercial unit at each address. Retain two homes per building, with respective floor areas of 4,375 and 4,750 square feet and reported land of 2,516 and 2,517 square feet. Retain the 2007 Assessor completion proxy; passed construction inspections occurred in early 2008. Their unchanged exact parcel identities support using current parcel locations with the construction-year ward map. This changes no land measurement.

Original Greenleaf permits specify nine apartments in each building; both passed in June 2007. The recorded correction retains eighteen homes in 2007, preserving the commercial source's whole-pair floor and land measurements. An explicitly reviewed partial residential copy can be retired only when the commercial replacement covers every source parcel. Automatic replacements still require identical parcel sets. The only additional active decision admitted by this change is Greenleaf. Diversey contains separate six- and three-apartment buildings, with commercial space excluded from residential counts. Winnebago's 2016 new-building permit and 2017 completed-building listing support five apartments in 2017; its repeated parcel reports should contribute one primary building measurement and the two reported lots. Bishop contains three separate four-flats; listing evidence supports a 2010 proxy at 15 and 17, while the located final inspection supports 2011 at 19. These recommendations are not all entered yet.

Polk's 2009 permits and first new assessments in 2010 conflict with the previously accepted 2008 Assessor/listing year. A revised 2010 proxy awaits Jacob. Lexington's twenty-four individual townhomes have 2012 permits, contradicting 2011; row inspections span 2013--2014, with some individual homes already sold in 2013. A common row-completion convention is a proposed proxy, not proof of every home's individual completion date, and awaits Jacob.

For Hope Manor II and four St. Edmund records, the remaining problem is measurement coverage. Hope's campus-wide reported land cannot be assigned to the reviewed subset of buildings without an unsupported allocation. St. Edmund's commercial counts total seventy-seven, conflicting with the developer's fifty-three-home development and building-specific permits. The recommendation is to withhold density for these five records until coherent reported measurements exist, preserving the evidence. That recommendation awaits Jacob; no campus totals are divided or allocated by hand.

Verification: construction Make and the targeted reconciliation report pass. An unchanged construction build requires no work. The fixed 268-case cohort has 258 closed cases and ten open cases. The candidate contains 13,762 unique projects. Compared with the pre-review candidate, only the two Michigan buildings were added and the three Lake/Greenleaf commercial records changed substantive fields; no unrelated years, counts, measurements or density permissions changed. All five retained project records have resolved locations. The 65th Place source is excluded. The paper input checksum is unchanged. No regression has been rerun.


\subsection*{Approved completion dates and density exclusions}
Jacob approved 2010 for the three Polk three-flats, row-completion proxies of 2013 and 2014 for Lexington's twenty-four separate townhomes, and 2010 for 15/17 Bishop and 2011 for 19 Bishop. He also approved withholding both density measures for Hope Manor II and the four St. Edmund records. Their source evidence remains recorded; no campus land or development-wide count is allocated to individual buildings.

During implementation, the historical address records identified Bishop parcel ending 052 as 15 North Bishop, not 19. The previous review had mislabeled this record. The correct current parcel for 19 ends in 050 and already supplies its own four-apartment assessment, with 7,175 square feet of floor area and 3,750 square feet of land. The obsolete 052 record is retired in favor of the completed four-home condominium at 15. The intended result is one building at each of 15, 17 and 19, using the approved years.

Diversey and Winnebago use recorded whole-building measurements in the existing project-selection script. The apartment count at 1136 Diversey excludes its shop. Winnebago uses the primary building area once across its two reported lots. These decisions are explicit inputs with evidence, rather than address-specific branches in code. Corrected years require additional official historical parcel queries; the existing acquisition task records and retrieves them. Parcel geometry supplies locations and ward-boundary distances, never replacement density denominators.

Verification is complete. All 268 cases in the fixed review cohort are closed and the current unresolved-case table has zero rows. The candidate contains 13,764 unique project records. Relative to the previous candidate, it adds 1130 Diversey, the completed 15 Bishop condominium, and Winnebago, while removing the obsolete 15 Bishop parcel. Thirty-three continuing records change approved years, apartment counts or density eligibility. All reviewed retained records have valid locations and distances using the approved construction years; no unrelated substantive fields change. The three added records are approximately 95, 188 and 387 feet from the relevant ward boundary, respectively.

Construction Make, source acquisition Make and the targeted audit pass. An unchanged construction build requires no work, and a deleted final-dataset report is regenerated from the saved data. The source refresh preserves every earlier downloaded record and geometry. The paper input remains unchanged and no regression has been rerun. These checks complete this review queue; they do not substitute for reconnecting the candidate to the analysis and running the paper.
